%iffalse
\let\negmedspace\undefined
\let\negthickspace\undefined
\documentclass[journal,12pt,twocolumn]{IEEEtran}
\usepackage{cite}
\usepackage{amsmath,amssymb,amsfonts,amsthm}
\usepackage{algorithmic}
\usepackage{graphicx}
\usepackage{textcomp}
\usepackage{xcolor}
\usepackage{txfonts}
\usepackage{listings}
\usepackage{enumitem}
\usepackage{mathtools}
\usepackage{gensymb}
\usepackage{comment}
\usepackage[breaklinks=true]{hyperref}
\usepackage{tkz-euclide} 
\usepackage{listings}
\usepackage{gvv}                                        
%\def\inputGnumericTable{}                                 
\usepackage[latin1]{inputenc}                                
\usepackage{color}                                            
\usepackage{array}                                            
\usepackage{longtable}                                       
\usepackage{calc}                                             
\usepackage{multirow}                                         
\usepackage{hhline}                                           
\usepackage{ifthen}                                           
\usepackage{lscape}
\usepackage{tabularx}
\usepackage{array}
\usepackage{float}


\newtheorem{theorem}{Theorem}[section]
\newtheorem{problem}{Problem}
\newtheorem{proposition}{Proposition}[section]
\newtheorem{lemma}{Lemma}[section]
\newtheorem{corollary}[theorem]{Corollary}
\newtheorem{example}{Example}[section]
\newtheorem{definition}[problem]{Definition}
\newcommand{\BEQA}{\begin{eqnarray}}
\newcommand{\EEQA}{\end{eqnarray}}
\newcommand{\define}{\stackrel{\triangle}{=}}
\theoremstyle{remark}
\newtheorem{rem}{Remark}

% Marks the beginning of the document
\begin{document}
\bibliographystyle{IEEEtran}
\vspace{3cm}

\title{Ch.5-Mathematical Induction \\
       I. Integer Value Correct Type}
\author{AI24BTECH11016-Jakkula Adishesh Balaji$^{*}$
}
\maketitle
\newpage
\bigskip


\renewcommand{\thefigure}{\theenumi}
\renewcommand{\thetable}{\theenumi}

1. The coefficients of three consecutive terms of $(1+x)^{n+5}$ are in the ratio 5:10:14. Then n=\\            (JEE Adv. 2013) \\ \\
2. Let m be the smallest positive integer such that the coefficient of $x^{2}$ in the expansion of $(1+x)^{2} + (1+x)^{3} + ... + (1+x)^{49} + (1+x)^{50} + (1+mx)^{50}$ is (3n+1) $\comb{51}{3}$ for some positive integer n. Then the value of n is \\ (JEE Adv. 2016) \\ \\
3. Let $\\ X= (\comb{10}{1})^{2}+2(\comb{10}{2})^{2} + 3(\comb{10}{3})^{2} + ... + 10(\comb{10}{10})^{2}$, where $\comb{10}{r} , r \in {1,2,...,10}$ denote binomial coefficients.\\ Then, the value of $\frac{1}{1430}$X is \\ (JEE Adv. 2018) \\ \\
4. Suppose \\
 \begin{vmatrix} $\sum\limits_{k=0}^{n}k$ & $\sum\limits_{k=0}^{n} k^{2} \comb{n}{k}$ \\ $\sum\limits_{k=0}^{n}\comb{n}{k} k$ & $\sum\limits_{k=0}^{n}\comb{n}{k} 3^{k}$
\end{vmatrix}
=0 \\	
holds for some positive integer n. The $\sum\limits_{k=0}^{n} \frac{\comb{n}{k}}{k+1}$ equals \, \, \, (JEE Adv. 2019)

\end{document}
